\documentclass[11pt,a4paper]{article}

\usepackage[margin=2.5cm]{geometry}
\usepackage{graphicx}
\usepackage{listings}
\usepackage{xcolor}
\usepackage{hyperref}
\usepackage{enumitem}

\lstdefinelanguage{SystemVerilog}{
    morekeywords={module,endmodule,input,output,wire,reg,logic,always,posedge,
        negedge,if,else,begin,end,case,endcase,assign,parameter,localparam,
        assert,assume,property,endproperty,disable,iff,initial,past,inside,
        default,ifdef,endif},
    sensitive=true,
    morecomment=[l]{//},
    morecomment=[s]{/*}{*/},
    morestring=[b]"
}

\lstset{
    basicstyle=\ttfamily\footnotesize,
    keywordstyle=\color{blue}\bfseries,
    commentstyle=\color{gray}\itshape,
    stringstyle=\color{red!70!black},
    breaklines=true,
    showstringspaces=false,
    frame=single,
    rulecolor=\color{black!30},
    language=SystemVerilog,
    tabsize=4,
    captionpos=b
}

\title{PV Lab 7 -- Formal Verification of the Packet FSM}
\author{Dan Joshua Hagen}
\date{\today}

\begin{document}

\maketitle

\noindent\textbf{Hours spent on this lab:} TBD

\section{FSM Graph}

\begin{figure}[h!]
    \centering
    \includegraphics[width=0.9\textwidth]{graph.png}
    \caption{Finite state machine of the \texttt{packet} module.}
    \label{fig:fsm}
\end{figure}

\section{Packet Formal Statements}

This section lists the \texttt{assume}/\texttt{assert} statements added to the
\texttt{packet.sv} module to formally verify it with EBMC. The full set is
located between the \texttt{`ifdef FORMAL} / \texttt{`endif} guards.

\subsection{Assumptions}

These constrain the environment so that the formal tool only explores legal
scenarios:

\begin{itemize}[leftmargin=*]
    \item \textbf{Start in reset.} The design is only meant to operate after a
        proper reset, so we force the very first cycle to be a reset cycle and
        require the state to enter \texttt{IDLE} on reset. This avoids
        spurious counter-examples that start from arbitrary register values.
    \item \textbf{\texttt{chk\_ok} and \texttt{chk\_fail} are mutually
        exclusive.} A real checksum check cannot report both \emph{pass} and
        \emph{fail} in the same cycle, so we tell the tool not to explore that
        input combination.
\end{itemize}

\begin{lstlisting}[caption={Assumptions on the environment.}]
initial assume(reset);
always @(posedge clk) if (reset) assume(state == IDLE);

assume property (@(posedge clk)
    !(chk_ok && chk_fail));
\end{lstlisting}

\subsection{Assertions}

\subsubsection{State transitions}

For every legal stimulus we assert that the FSM advances to the expected next
state. There is one assertion per transition in the FSM graph.

\textit{e.g.} from \texttt{HEADER}, asserting \texttt{hdr\_done} (and not
\texttt{abort}/\texttt{reset}) must move the FSM to \texttt{PAYLOAD}:

\begin{lstlisting}[caption={Example: HEADER $\rightarrow$ PAYLOAD on \texttt{hdr\_done}.}]
assert property (@(posedge clk) disable iff (reset)
    past_valid && !$past(reset || abort) && ($past(hdr_done) && $past(state) == HEADER)
        |-> state == PAYLOAD);
\end{lstlisting}

Analogous assertions cover \texttt{IDLE}$\rightarrow$\texttt{HEADER},
\texttt{PAYLOAD}$\rightarrow$\texttt{CHECKSUM},
\texttt{CHECKSUM}$\rightarrow$\texttt{DONE} (with \texttt{valid\_pkt}),
\texttt{CHECKSUM}$\rightarrow$\texttt{IDLE} (with \texttt{error\_pkt}), and
the automatic \texttt{DONE}$\rightarrow$\texttt{IDLE}.

\subsubsection{State holding}

When the triggering input is not asserted (and no reset/abort), the FSM must
remain in its current state. One assertion per such state.

\textit{e.g.} \texttt{HEADER} stays in \texttt{HEADER} while \texttt{hdr\_done}
is low:

\begin{lstlisting}[caption={Example: state holding in \texttt{HEADER}.}]
assert property (@(posedge clk) disable iff (reset)
    past_valid && !$past(reset || abort) && (!$past(hdr_done) && $past(state) == HEADER)
        |-> state == HEADER);
\end{lstlisting}

\subsubsection{Abort and reset go to \texttt{IDLE}}

From any state, asserting \texttt{abort} or \texttt{reset} must bring the FSM
back to \texttt{IDLE} on the next cycle.

\begin{lstlisting}[caption={Reset/abort dominate.}]
assert property (@(posedge clk)
    $past(reset || abort) |-> state == IDLE);
\end{lstlisting}

\subsubsection{Reset clears the output signals}

After a reset cycle, both \texttt{valid\_pkt} and \texttt{error\_pkt} must be
low.

\begin{lstlisting}[caption={Reset clears outputs.}]
assert property (@(posedge clk)
    $past(reset) |-> !valid_pkt && !error_pkt);
\end{lstlisting}

\subsubsection{Output signals are single-cycle pulses tied to the correct
transition}

\texttt{valid\_pkt} (resp. \texttt{error\_pkt}) may only be high in the cycle
immediately after a successful (resp. failed) checksum, and never together.
This catches the FSM raising a flag in the wrong state or holding it high.

\begin{lstlisting}[caption={Output flags are mutually exclusive and bound to the right transition.}]
assert property (@(posedge clk) disable iff (reset)
    !(valid_pkt && error_pkt));

assert property (@(posedge clk) disable iff (reset)
    valid_pkt |-> $past(state) == CHECKSUM && $past(chk_ok) && !$past(abort));

assert property (@(posedge clk) disable iff (reset)
    error_pkt |-> $past(state) == CHECKSUM && $past(chk_fail) && !$past(chk_ok) && !$past(abort));
\end{lstlisting}

\section{Tool Output -- Clean Design}

% TODO: paste EBMC output here once the tool is run.

\section{Tool Output -- Injected Bug}

% TODO: describe the bug, paste the counter-example, and comment on the waveform.

\end{document}
